\kbvsechead{(೬)\hspace{0.45em}ಶ್ಲೇಷಾಲಂಕಾರ :}

\indent “ಆ ಉದ್ಯಾನವು ಅಮರಾವತಿಯಂತೆ ಸುಮನಭರಿತವಾಗಿತ್ತು.”—ಈ ವಾಕ್ಯದಲ್ಲಿ ಬಳಸಿರುವ “ಸುಮನ” ಎಂಬ ಶಬ್ದಕ್ಕೆ ಎರಡು ಅರ್ಥಗಳುಂಟು ; ವರ್ಣಿಸಬೇಕಾದ ಉದ್ಯಾನದ ಪರವಾಗಿ ಹೇಳುವುದಾದರೆ, ಆ ಉದ್ಯಾ ನವು ಹೂವುಗಳಿಂದ ತುಂಬಿತ್ತು ಎಂದಾಗುವುದು ; ವರ್ಣಿಸುವ ವಿಷಯವಲ್ಲದ ಅಮರಾವತಿಯ ಪರವಾಗಿ ಹೇಳುವುದಾದರೆ, ಅಮರಾವತಿಯು ದೇವತೆಗಳಿಂದ ತುಂಬಿತ್ತು ಎಂದಾಗುವುದು. ಇಲ್ಲಿ ಶ್ಲೇಷೆಯಿರುವುದನ್ನು ಕಾಣಬಹುದು, {[}ಅಮರಾವತಿಯು ದೇವತೆಗಳಿಂದ ತುಂಬಿರುವಂತೆ, ಆ ಉದ್ಯಾನವು ಹೂವುಗಳಿಂದ ತುಂಬಿತ್ತು ಎನ್ನುವುದು ಒಟ್ಟಾರೆ ಅಭಿಪ್ರಾಯ{]}. ಹೀಗೆ, ಒಂದೇ ಶಬ್ದವು ವರ್ಣಿಸಬೇಕಾದ ವಿಷಯದ ಪರವಾಗಿ ಒಂದು ಅರ್ಥವನ್ನೂ ವರ್ಣನೆಯ ಅಗತ್ಯವಿರದ ವಿಷಯದ ಪರವಾಗಿ ಮತ್ತೊಂದು ಅರ್ಥವನ್ನು ಕೊಡುವುದಾದರೆ, ಅಂತಹ ಅರ್ಥಾಲಂಕಾರವು ಶ್ಲೇಷಾಲಂಕಾರ ವೆನಿಸುವುದು.\par

\vspace{7pt plus 6pt minus 1pt}\begin{center}{\kbvsecfont\fontsize{13pt}{16.5pt}\selectfont\color{kbvaccentdark} 2. ಶಬ್ದಾಲಂಕಾರ}\end{center}\vspace{3pt plus 3pt minus 1pt}

\indent ಶಬ್ದಾಲಂಕಾರದಲ್ಲಿ ಹಲವು ಪ್ರಕಾರಗಳುಂಟು, ಅವುಗಳು ಯಾವು ವೆಂಬುದನ್ನು ನೋಡೋಣ.\par

\kbvsechead{(೧)\hspace{0.45em}ವೃತ್ತಾನುಪ್ರಾಸ :}

\indent “ಕೊಳುಗುಳದ ಕೋಳ್ಗುದಿಯ ಕೋಲಾಹಲದ ಕೆಸರಿನಲಿ”-ಇದನ್ನು ಕರ್ಣಾಟಕ ಭಾರತ ಕಥಾಮಂಜರಿಯಿಂದ ಆರಿಸಿ ತೆಗೆದಿದೆ. ಇಲ್ಲಿ ಒಂದೇ ಅಕ್ಷರ (ಕ-ಕಾರ)ವು ನಾಲ್ಕು ಬಾರಿ ಪುನರುಕ್ತವಾಗಿ ವಾಕ್ಯದ ಸೊಬಗನ್ನು ಹೆಚ್ಚಿಸಿದೆ, ಈ ರೀತಿಯಾಗಿ, ಒಂದು ಅಥವಾ ಎರಡು ವ್ಯಂಜನಾಕ್ಷರಗಳು ಹಲವು ಬಾರಿ ಪುನರುಕ್ತವಾಗಿದ್ದರೆ ಅದನ್ನು ವೃತ್ತಾನುಪ್ರಾಸವೆನ್ನುವರು,\par
